\begin{table*}[t]
\centering
\small
\setlength{\tabcolsep}{4pt}
\begin{tabular}{l | rrrr | rrr | rr}
\toprule
 & \multicolumn{4}{c|}{Solve Rate (\%)}
 & \multicolumn{3}{c|}{Cost \& Time}
 & \multicolumn{2}{c}{Quality} \\
\cmidrule(lr){2-5} \cmidrule(lr){6-8} \cmidrule(l){9-10}
Model & Strict & Iso. & Core & Partial & \$/CKPT & Net \$ & Min/CKPT & Erosion & Verbosity \\
\midrule
Sonnet 4.5 & 5.4 & 16.1 & 34.4 & 15.0 & 1.49{\scriptsize$\pm$0.59} & 138.47 & 8.8{\scriptsize$\pm$3.1} & 0.682{\scriptsize$\pm$0.185} & 0.293{\scriptsize$\pm$0.126} \\
Sonnet 4.6 & 8.5 & 18.3 & 45.1 & 25.0 & 1.92{\scriptsize$\pm$1.51} & 157.38 & 20.0{\scriptsize$\pm$20.7} & 0.703{\scriptsize$\pm$0.251} & 0.313{\scriptsize$\pm$0.112} \\
Opus 4.5 & 10.9 & 17.4 & 44.6 & 25.0 & 2.64{\scriptsize$\pm$1.60} & 242.64 & 7.6{\scriptsize$\pm$4.2} & 0.710{\scriptsize$\pm$0.191} & 0.287{\scriptsize$\pm$0.083} \\
Opus 4.6 & \textbf{17.2} & 21.5 & \textbf{53.8} & \textbf{35.0} & 3.47{\scriptsize$\pm$3.40} & 322.97 & 14.4{\scriptsize$\pm$10.9} & 0.774{\scriptsize$\pm$0.132} & 0.346{\scriptsize$\pm$0.102} \\
\midrule
GPT~5.1 Codex Max & 10.8 & 17.2 & 38.7 & 25.0 & 2.86{\scriptsize$\pm$2.39} & 266.39 & 14.0{\scriptsize$\pm$9.5} & 0.642{\scriptsize$\pm$0.158} & 0.331{\scriptsize$\pm$0.095} \\
GPT~5.2 & 10.8 & 19.4 & 43.0 & 25.0 & 4.55{\scriptsize$\pm$6.34} & 422.72 & 15.0{\scriptsize$\pm$11.8} & 0.711{\scriptsize$\pm$0.212} & 0.358{\scriptsize$\pm$0.095} \\
GPT~5.2 Codex & 9.7 & 18.3 & 33.3 & 25.0 & 2.89{\scriptsize$\pm$3.12} & 268.78 & 14.6{\scriptsize$\pm$9.8} & 0.689{\scriptsize$\pm$0.209} & 0.388{\scriptsize$\pm$0.120} \\
GPT~5.3 Spark & 5.4 & 12.9 & 26.9 & 15.0 & \textbf{0.91{\scriptsize$\pm$2.17}} & \textbf{84.46} & \textbf{2.8{\scriptsize$\pm$6.1}} & 0.575{\scriptsize$\pm$0.282} & 0.352{\scriptsize$\pm$0.104} \\
GPT~5.3 Codex & 9.7 & \textbf{23.7} & 51.6 & 30.0 & 3.14{\scriptsize$\pm$3.11} & 292.38 & 7.5{\scriptsize$\pm$4.6} & 0.676{\scriptsize$\pm$0.156} & 0.356{\scriptsize$\pm$0.087} \\
GPT~5.4 & 11.8 & 20.4 & 48.4 & 30.0 & 3.27{\scriptsize$\pm$2.93} & 304.46 & 8.6{\scriptsize$\pm$5.3} & \textbf{0.515{\scriptsize$\pm$0.182}} & \textbf{0.286{\scriptsize$\pm$0.064}} \\
\midrule
GLM~4.7 & 4.3 & 9.7 & 32.3 & 15.0 & 1.61{\scriptsize$\pm$1.12} & 149.68 & 9.1{\scriptsize$\pm$5.3} & 0.664{\scriptsize$\pm$0.210} & 0.305{\scriptsize$\pm$0.073} \\
\bottomrule
\end{tabular}
\caption{Main \scb{} results for one predetermined harness version and one high-thinking, just-solve run per model. \textbf{Strict} requires all tests (including regression) to pass; \textbf{Iso.}\ excludes regression tests; \textbf{Core} counts only specification-demonstrated behavior; \textbf{Partial} is the fraction of problems with $\geq 1$ strictly solved checkpoint. Cost and time are per produced checkpoint. Erosion and verbosity are conditional on the agent producing a workspace. \textbf{Bold} marks the best value in each column.}
\label{tab:performance-overview}
\end{table*}

\autoref{tab:performance-overview} summarizes solve rates across \overallNumConfigs{} configurations and \overallNumModels{} models on \scb. No agent \emph{fully} solves any of the 20 problems: no run passes every test at every checkpoint end-to-end. \textbf{\overallBestModel{}} achieves the highest strict solve rate at \overallBestStrictSolveRate{}\%; isolated solve rates span \overallIsoSolveMin--\overallIsoSolveMax{}\%, core rates \overallCoreSolveMin--\overallCoreSolveMax{}\%.
\begin{figure}[t]
    \centering
    \includegraphics[width=\textwidth]{figures/overview_solve_cost.pdf}
    \caption{Solve rates and cost growth over problem progress. \emph{Left:} Agents pass all core tests \abFigCoreIsoRatioMin--\abFigCoreIsoRatioMax{}$\times$ more often than the full checkpoint suite, and strict solve rates, which include regression, collapse to \abFigStrictSolveFinal\% by the final checkpoint. \emph{Right:} Mean cost per checkpoint grows \abFigCostGrowth{}$\times$ from the first to the last progress bin. See \autoref{sec:test-pass-rates} for a breakdown of continuous pass rates by test type.}
    \label{fig:solve-cost}
\end{figure}

As checkpoints advance, the gap between core and isolated pass rates widens from \abFigCoreIsoRatioMin{}$\times$ to \abFigCoreIsoRatioMax{}$\times$ (\autoref{fig:solve-cost}). The newly introduced tests at each checkpoint are harder to satisfy than the original core suite, and error-handling tests account for most of the decline while core and functionality pass rates remain comparatively stable (\autoref{sec:test-pass-rates}). Cost grows \abFigCostGrowth{}$\times$ over the same span, but the additional spending does not improve correctness.

The remainder of these results examines how quality issues accumulate, how agent behaviors differ from human developers, and the impact of prompt instructions on quality degradation.


